% ===== Nervous breakdowns: "Emergency shutdown" (~16 min), then the break =====
\partaccent{peach}
\topicopener{battery}{Սթրես}{Վթարային անջատում}{Նյարդային պոռթկումներ, և ինչպես են դրանք առաջանում}{panic}

\begin{frame}{Մարմինը սահմաններ ունի}
  \begin{columns}[c]
    \begin{column}{0.5\textwidth}
      \begin{popup}{խորհուրդներ.txt}
        \small
        \begin{itemize}\setlength{\itemsep}{0.3mm}
          \item Առանց վերարկուի դուրս ես եկել՝ կմրսես
          \item Արագ սնունդը կփչացնի սիրտդ
          \item Ծխելը քայքայում է թոքերը
          \item Ծանրությունը բարձրացրու ծնկներով
          \item Մատներդ մի ճտտացրու, արթրիտ կունենաս
        \end{itemize}
      \end{popup}
      {\scriptsize\color{plum!70!cream} Բըրնեթ. «Սրանց ճշգրտությունը շատ տարբեր է»}
    \end{column}
    \begin{column}{0.44\textwidth}
      \onslide<2->{%
      \begin{card}[colback=lavender!35!cardwhite]
        Բայց գաղափարը ճիշտ է. մեր արարքներն ազդում են առողջության վրա։ Իսկ ի՞նչն է նույնից պաշտպանում ուղեղին։ \textbf{Ոչինչ}։
      \end{card}}
      \onslide<3->{%
      \begin{card}[colback=brainpink!30!cardwhite]
        \centering Այսօր ուղեղի ամենամեծ վտանգը՝ հին ու բարի \textbf{սթրեսն} է
      \end{card}}
    \end{column}
  \end{columns}
\end{frame}

\begin{frame}{Սթրեսի առանցքը}
  \begin{columns}[c]
    \begin{column}{0.4\textwidth}
      \centering
      \begin{tikzpicture}[ax/.style={flow=3.6cm, font=\small}]
        \node[ax, fill=lavender!40!cardwhite] (h) at (0,2.1) {\emoji{brain}\ \textbf{Հիպոթալամուս}};
        \node[ax, fill=lavender!40!cardwhite] (p) at (0,1.05) {\textbf{Հիպոֆիզ} (pituitary)};
        \node[ax, fill=peach!45!cardwhite] (a) at (0,0) {\textbf{Մակերիկամներ} (adrenal)};
        \node[ax, fill=brainpink!40!cardwhite, visible on=<2->] (c) at (0,-1.15) {\emoji{fire}\ ադրենալին ու \textbf{կորտիզոլ}};
        \draw[arr] (h) -- (p); \draw[arr] (p) -- (a);
        \draw[arr, visible on=<2->] (a) -- (c);
      \end{tikzpicture}\par\smallskip
      {\scriptsize HPA առանցք (hypothalamic-pituitary-adrenal)}
    \end{column}
    \begin{column}{0.56\textwidth}
      \begin{itemize}
        \item Սթրեսը միացնում է «կռվիր կամ փախիր» ռեակցիան (հիշո՞ւմ ես առաջին վեբինարից)
        \item<2-> Կորտիզոլը՝ «սթրեսի հորմոնը», ազդում է ամբողջ մարմնի ու ուղեղի վրա
        \item<3-> Երկար սթրեսից մարդը դառնում է լարված, դյուրագրգիռ, հյուծված, չի կարողանում պարզ մտածել
        \item<4-> Ասում ենք՝ «պոռթկման եզրին է»
      \end{itemize}
    \end{column}
  \end{columns}
\end{frame}

\begin{frame}{«Նյարդային պոռթկում»}
  \begin{columns}[T]
    \begin{column}{0.5\textwidth}
      \begin{card}[colback=butter!45!cream]
        {\blackfont Պաշտոնական տերմին չէ}\par\smallskip
        Նյարդերը բառացիորեն չեն «պոռթկում»։ Մարդն այլևս չի դիմանում սթրեսին ու կարծես «կոտրվում» է. փակվում է, քաշվում է, չի կարողանում նորմալ գործել
      \end{card}
    \end{column}
    \begin{column}{0.44\textwidth}
      \onslide<2->{%
      {\blackfont\small Ամեն մեկի մոտ՝ տարբեր}\par\smallskip
      \chip[lavender]{pensive}{խոր դեպրեսիա}\par\smallskip
      \chip[peach]{fearful}{խուճապի նոպաներ}\par\smallskip
      \chip[brainpink]{speakinghead}{երբեմն՝ հալյուցինացիաներ}}
    \end{column}
  \end{columns}
  \vspace{3mm}
  \onslide<3->{%
  \begin{card}[colback=mint!30!cardwhite]
    \emoji{bandage}\ Ոմանք պոռթկումը համարում են ուղեղի \textbf{պաշտպանական մեխանիզմ}. ինչպես ֆիզիոթերապիան. տհաճ է ու հյուծող, բայց ավելի լավ է, քան չանելը
  \end{card}}
\end{frame}

\begin{frame}{Սթրեսը վերահսկողության կորուստ է}
  \begin{columns}[c]
    \begin{column}{0.46\textwidth}
      \begin{itemize}
        \item Սթրես առաջացնողը (stressor) \textbf{նվազեցնում է վերահսկողությունը}
        \item Երբ թվում է, թե վերահսկում ենք իրավիճակը, ապահով ենք զգում
        \item Ու կարևոր չէ, թե իրականում որքան վերահսկողություն ունենք
      \end{itemize}
      \onslide<2->{\centering\bigemoji[1cm]{earth}\ \bigemoji[1cm]{coffee}\par}
    \end{column}
    \begin{column}{0.5\textwidth}
      \onslide<2->{%
      \begin{bookquote}
        Տեխնիկապես ամեն մարդ ածխածնի անիմաստ պարկ է, որ կառչել է միջուկային կրակի շուրջ անտարբեր դատարկության մեջ սլացող քարից։ Բայց եթե կարող ես սուրճիդ մեջ սոյայի կաթ պահանջել ու ստանալ, դա \textbf{շոշափելի վերահսկողություն} է։
      \end{bookquote}}
    \end{column}
  \end{columns}
\end{frame}

\begin{frame}{Երբ ոչինչ չես կարող անել}
  \begin{columns}[c]
    \begin{column}{0.4\textwidth}
      \centering
      \begin{meme}[4.35cm]{panik_kalm_panik.png}
        \memecap[text width=1.8cm, font=\blackfont\scriptsize]{0.245}{0.83}{Անձրև է գալիս}
        \memecap[text width=1.8cm, font=\blackfont\scriptsize]{0.245}{0.495}{Անձրևանոց ունեմ}
        \memecap[text width=1.8cm, font=\blackfont\scriptsize]{0.245}{0.165}{Անձրևանոցը տանն է, բանալիներն էլ}
      \end{meme}
    \end{column}
    \begin{column}{0.56\textwidth}
      \begin{itemize}
        \item Սթրեսն ավելի մեծ է, երբ \textbf{ընտրություն չկա}
        \item Գլխացավի կամ մրսածության համար դեղ կա
        \item Իսկ քրոնիկ հիվանդությունը մշտական, անխուսափելի տհաճություն է, ու դա շատ սթրեսային է
      \end{itemize}
    \end{column}
  \end{columns}
\end{frame}

\begin{frame}{Հոգնածություն ու անկանխատեսելիություն}
  \begin{columns}[T]
    \begin{column}{0.47\textwidth}
      \begin{card}[colback=peach!35!cardwhite]
        \emoji[1.3em]{batterylow}\ {\blackfont Հոգնեցնում է}\par\smallskip
        Քնել ես մնացել ու վազում ես գնացքի հետևից։ Սթրեսը էներգիա է ծախսում, իսկ հոգնածությունը նոր սթրես է
      \end{card}
    \end{column}
    \begin{column}{0.47\textwidth}
      \onslide<2->{%
      \begin{card}[colback=lavender!40!cardwhite]
        \emoji[1.3em]{dice}\ {\blackfont Անկանխատեսելի է}\par\smallskip
        Էպիլեպսիայի նոպան կարող է գալ ցանկացած պահի։ Կամ զուգընկեր, որ կարող է պայթել սուրճի բանկան սխալ դարակում դնելու համար
      \end{card}}
    \end{column}
  \end{columns}
  \vspace{1mm}
  \onslide<3->{%
  \begin{card}[colback=mint!30!cardwhite]
    \emoji{couch}\ Սովորաբար հաղթահարում ենք. պարասիմպաթիկ համակարգը («բազմոցի տղան» առաջին վեբինարից) հանգստացնում է մեզ
  \end{card}}
\end{frame}

\begin{frame}{Հոլմս-Ռահեի սթրեսի սանդղակը}
  \begin{columns}[c]
    \begin{column}{0.36\textwidth}
      {\small 1967. ամեն իրադարձություն ստացավ «կյանքի փոփոխության միավորներ» (LCU)}\par\medskip
      \begin{poll}
        Ո՞րն է ավելի սթրեսային. \textbf{ամուսնալուծությո՞ւնը}, թե՞ \textbf{բանտը}։
      \end{poll}
      \onslide<3->{{\small Բըրնեթ. «Մի տեսակ ռոմանտիկ է» \emoji{heart}}}
    \end{column}
    \begin{column}{0.62\textwidth}
      \centering
      \onslide<2->{%
      % Holmes & Rahe (1967) Social Readjustment Rating Scale, life change units (LCU).
% Blue/red = values in the book (red = the divorce vs prison poll); orange = from the same scale, not in the book.
% ytick lists every coordinate: with one \addplot per colour, ytick=data would only take the first plot's.
\begin{tikzpicture}
  \begin{axis}[
      xbar, width=6.4cm, height=4.9cm, bar width=0.3cm, bar shift=0pt,
      xmin=0, xmax=122, enlarge y limits=0.07,
      symbolic y coords={xmas, inlaws, fired, marriage, injury, prison, divorce, spouse},
      ytick={xmas, inlaws, fired, marriage, injury, prison, divorce, spouse},
      yticklabels={Սուրբ Ծնունդ, Վեճ խնամիների հետ, Աշխատանքից ազատում, Ամուսնություն, Անձնական վնասվածք, Բանտարկություն, Ամուսնալուծություն, Կողակցի մահ},
      tick label style={font=\scriptsize, text=plum}, axis line style={plum},
      axis x line*=bottom, axis y line*=left, xtick=\empty,
      nodes near coords, nodes near coords style={font=\blackfont\scriptsize, text=plum, anchor=west},
      point meta=x, clip=false,
    ]
    \addplot[fill=amBlue, draw=none] coordinates {(100,spouse) (53,injury) (47,fired) (29,inlaws)};
    \addplot[fill=amRed, draw=none] coordinates {(73,divorce) (63,prison)};
    \addplot[fill=amOrange, draw=none] coordinates {(50,marriage) (12,xmas)};
  \end{axis}
\end{tikzpicture}
\par
      {\scriptsize\color{amOrange!80!black}\textbf{նարնջագույնը}}{\scriptsize\ նույն սանդղակից է, գրքում չկա}}
    \end{column}
  \end{columns}
  \only<2->{\source{Holmes \& Rahe, J. Psychosom. Res., 1967}}
\end{frame}

\begin{frame}{Սուր սթրես}
  \begin{columns}[c]
    \begin{column}{0.54\textwidth}
      \begin{itemize}
        \item Վթար, բռնություն, մեծ ողբերգություն. մեկ դեպք, որից սթրեսը դառնում է \textbf{անտանելի}
        \item<2-> Սթրեսի հորմոնները կարճ ժամանակով \textbf{ուժեղացնում են հիշողությունը}. «ֆլեշ» հիշողություններ (flashbulb memories)
        \item<3-> Օգտակար է. վտանգը հիշում ես, որ էլ չկրկնես։ Բայց ծայրահեղ դեպքում հիշողությունն այնքան վառ է, որ մարդն այն \textbf{նորից ու նորից} է ապրում
      \end{itemize}
    \end{column}
    \begin{column}{0.4\textwidth}
      \onslide<3->{%
      \begin{card}[colback=butter!45!cream]
        \emoji[1.3em]{sun}\ Ինչպես շատ պայծառ լույսը, որ «այրվում» է աչքիդ մեջ։ Միայն թե չի անցնում
      \end{card}}
      \onslide<4->{%
      \begin{bookquote}
        Ուղեղի համակարգը, որ պետք է կանխի տրավմայի կրկնությունը, ինքն է այն կրկնում։
      \end{bookquote}}
    \end{column}
  \end{columns}
\end{frame}

\begin{frame}{Անջատվել ամեն ինչից}
  \begin{columns}[c]
    \begin{column}{0.48\textwidth}
      \begin{popup}{ուղեղ.exe}
        {\blackfont Կյանքը չափազանց սթրեսային է}\par\smallskip
        Անցնում եմ \textbf{սպասման ռեժիմի}... \emoji{zzz}
      \end{popup}
    \end{column}
    \begin{column}{0.48\textwidth}
      \begin{itemize}
        \item Մշտական վառ հիշողություններից մարդը կարող է «բթանալ» (numbing) կամ կտրվել իրականությունից (dissociation)
        \item<2-> Կարճ ժամանակով օգնում է, բայց երկար ժամանակով վնասում է մտածողությանն ու վարքին
        \item<3-> Ամենահայտնի հետևանքը՝ \textbf{հետտրավմատիկ սթրեսային խանգարումը} (PTSD)
      \end{itemize}
    \end{column}
  \end{columns}
\end{frame}

\begin{frame}{Քրոնիկ սթրես}
  \begin{columns}[c]
    \begin{column}{0.4\textwidth}
      \begin{itemize}
        \item Խնամել հիվանդ հարազատին, բռնակալ պետ, անվերջ դեդլայններ, պարտքեր
        \item<2-> Պարասիմպաթիկ համակարգը \textbf{հյուծվում} է, ու սթրեսը դառնում է «նորմա»
        \item<3-> Մարդը մշտապես լարված է, ցրված, ցնցվում է ամեն ձայնից
      \end{itemize}
    \end{column}
    \begin{column}{0.58\textwidth}
      \centering
      % Schematic, not data: after one stressful event the body is back to normal in ~20-60 min (book: Ulrich et al., 1991);
% with chronic stressors new spikes arrive before it recovers, so "stressed" becomes the new normal. Slide 2 adds the red curve.
\begin{tikzpicture}
  \begin{axis}[
      width=7.4cm, height=4.6cm,
      xmin=0, xmax=4, ymin=0, ymax=2.2,
      xtick={0,1,2,3,4}, ytick=\empty,
      xlabel={ժամեր}, ylabel={սթրեսի հորմոններ},
      label style={font=\scriptsize, text=plum}, tick label style={font=\scriptsize, text=plum},
      axis line style={plum}, axis x line*=bottom, axis y line*=left, clip=false,
    ]
    \draw[plum!40, dashed, line width=0.8pt] (axis cs:0,0.25) -- (axis cs:4,0.25);
    \node[anchor=north, font=\scriptsize, text=plum!70!cream, inner sep=1pt] at (axis cs:1.0,0.22) {նորմա};
    \addplot[amBlue, line width=1.6pt, domain=0:4, samples=120] {0.25 + (x>0.3)*0.9*exp(-(x-0.3)/0.22)};
    \node[anchor=south west, align=left, font=\blackfont\scriptsize, text=amBlue, inner sep=1pt] at (axis cs:0.4,1.2) {մեկ դեպք.\\ 20-60 րոպեում անցնում է};
    \addplot[amRed, line width=1.6pt, domain=0:4, samples=200, visible on=<2->]
      {0.25 + (x>1.5)*0.55*exp(-(x-1.5)/0.7) + (x>2.0)*0.55*exp(-(x-2.0)/0.7) + (x>2.5)*0.55*exp(-(x-2.5)/0.7)
             + (x>3.0)*0.55*exp(-(x-3.0)/0.7) + (x>3.5)*0.55*exp(-(x-3.5)/0.7)};
    \node[anchor=north east, align=right, font=\blackfont\scriptsize, text=amRed, inner sep=1pt, visible on=<2->] at (axis cs:4,2.2) {քրոնիկ սթրես.\\ չի հասցնում իջնել};
  \end{axis}
\end{tikzpicture}
\par
      {\scriptsize սխեմատիկ}
    \end{column}
  \end{columns}
\end{frame}

\begin{frame}{Մի քիչ սթրեսն օգնում է}
  \begin{columns}[c]
    \begin{column}{0.5\textwidth}
      \centering
      % Schematic inverted U of the Yerkes-Dodson law (the book's footnote): some pressure helps, too much hurts.
\begin{tikzpicture}
  \begin{axis}[
      width=6.6cm, height=4.6cm,
      xmin=0, xmax=1, ymin=0, ymax=1.25,
      xtick=\empty, ytick=\empty,
      xlabel={սթրես, ճնշում}, ylabel={արդյունք},
      label style={font=\scriptsize, text=plum},
      axis line style={plum}, axis x line*=bottom, axis y line*=left, clip=false,
    ]
    \addplot[amBlue, line width=1.8pt, domain=0:1, samples=80] {0.08 + 0.92*exp(-((x-0.45)/0.22)^2)};
    \node[font=\scriptsize, text=plum, align=center] at (axis cs:0.17,0.62) {ձանձրույթ\\ \emoji{yawn}};
    \node[font=\blackfont\scriptsize, text=mintdark, align=center] at (axis cs:0.45,1.15) {դեդլայնի «կախարդանք»};
    \node[font=\scriptsize, text=pinkdark, align=center, visible on=<2->] at (axis cs:0.88,0.55) {ամեն ինչ\\ քանդվում է\\ \emoji{exploding}};
  \end{axis}
\end{tikzpicture}

    \end{column}
    \begin{column}{0.46\textwidth}
      \begin{poll}
        Ձեռք բարձրացրու, եթե դեդլայնից առաջ ավելի լավ ես աշխատում։
      \end{poll}
      \begin{itemize}
        \item<2-> \textbf{Յերկս-Դոդսոնի օրենքը} (1908). մինչև որոշ կետ սթրեսն օգնում է, հետո՝ խանգարում
        \item<3-> Ինչպես աղը. քիչը համեղացնում է, շատը՝ փչացնում \emoji{salt}
      \end{itemize}
    \end{column}
  \end{columns}
\end{frame}

\begin{frame}{Յերկս-Դոդսոնը՝ գործատուի աչքերով}
  \centering
  \begin{meme}[4.6cm]{boardroom.jpg}
    \memecap[text width=3.0cm, font=\blackfont\scriptsize]{0.6}{0.945}{Ինչպե՞ս բարձրացնենք արդյունավետությունը}
    \memecap[text width=1.15cm, font=\blackfont\scriptsize]{0.15}{0.6}{Ավելի շատ դեդլայն}
    \memecap[text width=1.15cm, font=\blackfont\scriptsize]{0.41}{0.585}{Ավելի շատ ճնշում}
    \memecap[text width=1.75cm, font=\blackfont\scriptsize]{0.76}{0.56}{Շատ սթրեսն իջեցնում է արդյունքը}
  \end{meme}
\end{frame}

\begin{frame}{Սթրեսի ցիկլ}
  \centering
  % Burnett's "stress cycle" with his new-boss example. Steps reveal on slides 1-5, the loop closes on slide 6.
\begin{tikzpicture}[st/.style={flow=3.15cm, font=\small, inner sep=3pt}]
  \node[st, fill=butter!55!cardwhite] (n1) at (0,1.75) {\emoji{boss}\ Նոր պետ, անհնար ծավալ};
  \node[st, fill=peach!45!cardwhite, visible on=<2->] (n2) at (3.15,0.5) {\emoji{hourglass}\ Երկար ժամեր, քրոնիկ սթրես};
  \node[st, fill=peach!65!cardwhite, visible on=<3->] (n3) at (1.95,-1.45) {\emoji{burger}\emoji{beer}\ Արագ սնունդ ու ալկոհոլ՝ «հանգստանալու»};
  \node[st, fill=brainpink!40!cardwhite, visible on=<4->] (n4) at (-1.95,-1.45) {\emoji{weary}\ Առողջությունն ու տրամադրությունը վատանում են};
  \node[st, fill=brainpink!60!cardwhite, visible on=<5->] (n5) at (-3.15,0.5) {\emoji{batterylow}\ Ավելի խոցելի ես նոր սթրեսի համար};
  \draw[arr, visible on=<2->] (n1.east) to[bend left=25] (n2.north);
  \draw[arr, visible on=<3->] (n2.south) to[bend left=20] (n3.east);
  \draw[arr, visible on=<4->] (n3.west) to[bend left=15] (n4.east);
  \draw[arr, visible on=<5->] (n4.west) to[bend left=20] (n5.south);
  \draw[arr=pinkdark, visible on=<6->] (n5.north) to[bend left=25] (n1.west);
  \node[font=\blackfont\small, text=pinkdark, visible on=<6->] at (0,0.1) {\emoji{cyclone}\ ու նորից};
\end{tikzpicture}
\par\medskip
  \onslide<6->{\small Արտաքին «թեթևացումը» հաճախ նոր սթրես է բերում։\par}
\end{frame}

\begin{frame}{Ապահովիչը թռավ}
  \begin{columns}[c]
    \begin{column}{0.5\textwidth}
      \begin{itemize}
        \item Ցիկլից դուրս գալու ձևեր կան. ծանրաբեռնվածությունը փոխել, առողջ ապրելակերպ, թերապիա
        \item<2-> Բայց հաճախ դա չի լինում, սթրեսը կուտակվում է, ու մի պահ \textbf{շեմն անցնում է}
        \item<3-> Ուղեղն ասես հանձնվում է ու «անջատում է ամեն ինչ»
      \end{itemize}
    \end{column}
    \begin{column}{0.46\textwidth}
      \onslide<3->{%
      \begin{popup}{ուղեղ.exe}
        \centering\bigemoji[1cm]{plug}\ \bigemoji[1cm]{lightning}\par\smallskip
        {\blackfont Վթարային անջատում}\par
        {\small Ինչպես ապահովիչը, որ անջատում է հոսանքը, նախքան ամբողջ համակարգը կայրվի}
      \end{popup}}
      \onslide<4->{\small Ոմանց կարծիքով՝ պոռթկումը կանխում է \textbf{երկարատև} վնասը։}
    \end{column}
  \end{columns}
\end{frame}

\begin{frame}{Ո՞ւմ մոտ է շեմն ավելի ցածր}
  \begin{columns}[c]
    \begin{column}{0.44\textwidth}
      \centering
      % Diathesis-stress model, schematic: vulnerability + stress reaching the same breaking threshold.
% A needs a lot of stress (slide 2), B, more vulnerable, needs little (slide 3).
\begin{tikzpicture}
  \fill[amBlue] (0,0) rectangle (1.5,0.9);
  \fill[amOrange, visible on=<2->] (0,0.9) rectangle (1.5,3.6);
  \fill[amBlue] (2.6,0) rectangle (4.1,2.75);
  \fill[amOrange, visible on=<3->] (2.6,2.75) rectangle (4.1,3.6);
  \draw[plum, line width=0.8pt] (-0.3,0) -- (4.5,0);
  \draw[amRed, dashed, line width=1.4pt] (-0.3,3.6) -- (4.5,3.6) node[right, font=\blackfont\scriptsize, text=amRed] {շեմ};
  \node[font=\blackfont\small, text=plum] at (0.75,-0.3) {Ա};
  \node[font=\blackfont\small, text=plum] at (3.35,-0.3) {Բ};
  % legend
  \fill[amBlue] (0,4.4) rectangle (0.3,4.65);
  \node[anchor=west, font=\scriptsize, text=plum] at (0.35,4.53) {խոցելիություն (դիաթեզ)};
  \fill[amOrange] (0,3.95) rectangle (0.3,4.2);
  \node[anchor=west, font=\scriptsize, text=plum] at (0.35,4.08) {սթրես};
\end{tikzpicture}

    \end{column}
    \begin{column}{0.52\textwidth}
      {\blackfont\small Դիաթեզ-սթրես մոդել} {\small (diathesis՝ խոցելիություն)}\par\smallskip
      \begin{itemize}
        \item<2-> Որքան մեծ է խոցելիությունը, այնքան \textbf{քիչ} սթրես է պետք
        \item<3-> Ավելի խոցելի են դժվար կյանքով ապրողները, անհանգստության հակված մարդիկ
        \item<4-> Նույնիսկ շատ ինքնավստահները. վերահսկողության կորուստը քանդում է նրանց ինքնապատկերը
      \end{itemize}
    \end{column}
  \end{columns}
\end{frame}

\begin{frame}{Ինչ է լինում պոռթկման ժամանակ ու հետո}
  \begin{columns}[c]
    \begin{column}{0.56\textwidth}
      \begin{itemize}
        \item Եթե արդեն կա հակում դեպրեսիայի կամ անհանգստության, սթրեսը կարող է այն «արթնացնել»
        \item<2-> Կորտիզոլը գրգռում է \textbf{դոֆամինի} համակարգը, որը կապված է փսիխոզի հետ
        \item<3-> Սովորաբար պոռթկումը \textbf{կարճատև} է. բժշկի օգնությամբ կամ հարկադիր դադարից հետո մարդիկ վերադառնում են կյանքին
        \item<4-> Բայց ոմանց մոտ սթրեսի նկատմամբ զգայունությունը \textbf{մնում է}
      \end{itemize}
    \end{column}
    \begin{column}{0.38\textwidth}
      \begin{card}[colback=peach!35!cardwhite]
        \emoji[1.3em]{books}\ Գիրքն ընկավ ոտքիդ մատին. ցավում է։ Իսկ եթե մատն արդեն \textbf{կոտրված} էր, շատ ավելի է ցավում
      \end{card}
    \end{column}
  \end{columns}
\end{frame}

\begin{frame}{Շնորհակալ լինե՞լ ուղեղին}
  \begin{columns}[c]
    \begin{column}{0.6\textwidth}
      \begin{bookquote}
        Ուղեղին շնորհակալ լինել, որ պոռթկմամբ սահմանափակում է սթրեսի վնասը, նույնն է, թե շնորհակալություն հայտնես մեկին, որ օգնեց տանդ հրդեհը մարել, երբ հենց ինքն էր տապակը միացրած թողել։
      \end{bookquote}
      \vspace{2mm}
      {\small Ուղեղի սթրեսը հաղթահարելու ձևերը ստեղծվել են ուրիշ աշխարհի համար ու ժամանակակից կյանքին միշտ չէ, որ համապատասխանում են։}
    \end{column}
    \begin{column}{0.34\textwidth}
      \centering\mascot[2.8cm]{smug}\par
      \bigemoji[1cm]{fire}
    \end{column}
  \end{columns}
\end{frame}

% ---------------------------------------------------------------- break
\begin{frame}[plain]
  \begin{tikzpicture}[remember picture, overlay]
    \fill[butter!60!cream] (current page.south west) rectangle (current page.north east);
    \node at ([yshift=8mm]current page.center) {\mascot[4cm]{sleepy}};
    \node[anchor=north, font=\titlefont\fontsize{30}{32}\selectfont, text=plum] at ([yshift=-17mm]current page.center) {Ընդմիջում};
    \node[anchor=north, font=\mediumfont\large, text=plum] at ([yshift=-31mm]current page.center)
      {5 րոպե \emoji{coffee}\quad Պլանային անջատում, ոչ թե վթարային};
  \end{tikzpicture}
\end{frame}
